\documentclass[a4paper,12pt]{article}
\usepackage{color}
\usepackage{graphicx, gensymb}
\usepackage{amsfonts, cancel, amssymb}
\usepackage{amsmath, amsthm, array, esvect, esint}
\usepackage{typearea}
\usepackage[a4paper,left=2cm,right=2cm,top=2cm,bottom=3cm,bindingoffset=5mm]{geometry}
\newcommand\numberthis{\addtocounter{equation}{1}\tag{\theequation}}
\newcommand{\bra}{}
\newcommand{\kom}{}
\newcommand{\brak}{}
\newcommand{\braket}{}
\newcommand{\intx}{}
\newcommand{\intr}{}
\newcommand{\kla}{}
\newcommand{\klam}{}
\newcommand{\klammer}{}
\newcommand{\oper}{}
\newcommand{\tecxt}{}
\newcommand{\Bra}{}
\newcommand{\Ket}{}

\begin{document}

\title{04 Lennard-Jones-Potential}
\author{Joachim Schwardt}
\date{\today}
\maketitle

\section*{Task 4.1: Binding energies of a neon-crystal}
\subsection*{a)}Let $U(r)=4\epsilon\klam\left[ {\kla\left( {\frac{\sigma}{r}} \right)^{12}-\kla\left( {\frac{\sigma}{r}} \right)^6} \right]$. 
\begin{align*}
0&\overset{!}{=}U'(r)\propto \frac{12}{r}\kla\left( {\frac{\sigma}{r}} \right)^{12}-\frac{6}{r}\kla\left( {\frac{\sigma}{r}} \right)^6&\Rightarrow &&\frac{\sigma}{r}&=\frac{1}{\sqrt[6]{2}}
\end{align*}
Therefor the equilibrium distance is $r_0=\sigma\sqrt[6]{2}$ and we get $U(r_0)=-\epsilon$.
\subsection*{b)}The full potential is $u(r):=\frac{U(r)}{N/2}=4\epsilon\klam\left[ {A_{12}\kla\left( {\frac{\sigma}{r}} \right)^{12}-A_6\kla\left( {\frac{\sigma}{r}} \right)^6} \right]$. The new equilibrium distance then becomes $r_0=\sigma\sqrt[6]{\frac{2A_{12}}{A_6}}$. 
\begin{align*}
\frac{u^\tecxt\text{bcc}}{\epsilon}&=4\klam\left[ {A_{12}^\tecxt\text{bcc}\kla\left( {\frac{\sigma}{r_0}} \right)^{12}-A_6^\tecxt\text{bcc}\kla\left( {\frac{\sigma}{r_0}} \right)^6} \right]=-\frac{(A_6^\tecxt\text{bcc})^2}{A_{12}^\tecxt\text{bcc}}=-16.473\\
\frac{u^\tecxt\text{hcp}}{\epsilon}&=4\klam\left[ {A_{12}^\tecxt\text{bcc}\kla\left( {\frac{\sigma}{r_0}} \right)^{12}-A_6^\tecxt\text{bcc}\kla\left( {\frac{\sigma}{r_0}} \right)^6} \right]=-\frac{(A_6^\tecxt\text{hcp})^2}{A_{12}^\tecxt\text{hcp}}=-17.222\\
\frac{u^\tecxt\text{fcc}}{\epsilon}&=4\klam\left[ {A_{12}^\tecxt\text{bcc}\kla\left( {\frac{\sigma}{r_0}} \right)^{12}-A_6^\tecxt\text{bcc}\kla\left( {\frac{\sigma}{r_0}} \right)^6} \right]=-\frac{(A_6^\tecxt\text{fcc})^2}{A_{12}^\tecxt\text{fcc}}=-17.220
\end{align*} 
%\begin{align*}
%A_k^\tecxt\text{fcc}&=\sum_{(x,y,z)\in\mathbb{Z}^3\setminus\{0\}}\frac{1}{\sqrt{x^2+y^2+z^2}^k}=\sum_{j\ge 1}\frac{1}{\sqrt{j}^k}\sum_{x^2+y^2+z^2=j}=\sum_{j\ge 1}\frac{4j^2+2}{j^{k/2}}
%\end{align*}
%\begin{align*}
%I&=\intx\int_{-a}^{b}\tanh(x+a)\,\mathrm{d}x-\intx\int_{a}^{b}\tanh(x-a)\,\mathrm{d}x=\\
%&=\log\frac{\cosh(a+b)}{\cosh(a-a)}-\log\frac{\cosh(a-b)}{\cosh(a-a)}=\log\frac{\cosh(a+b)}{\cosh(a-b)}=\\
%&=\log\frac{e^{a}+e^{-a-2b}}{e^{a-2b}+e^{-a}}\rightarrow\log\frac{e^a}{e^{-a}}=2a
%\end{align*}

\end{document}